% Student 4 Individual Report
% Project: Diffusion Policy Imitation Learning for Dual-Arm Manipulation
\documentclass[11pt]{article}

\usepackage[utf8]{inputenc}
\usepackage{times}
\usepackage{graphicx}
\usepackage{amsmath}
\usepackage{amssymb}
\usepackage{booktabs}
\usepackage{multirow}
\usepackage{subcaption}
\usepackage{xcolor}
\usepackage{float}
\usepackage{geometry}
\usepackage{hyperref}
\usepackage{listings}
\geometry{a4paper, margin=1in}

% Code listing style
\definecolor{codegray}{rgb}{0.95,0.95,0.95}
\definecolor{codecomment}{rgb}{0.0,0.5,0.0}
\definecolor{codekeyword}{rgb}{0.0,0.0,0.7}
\definecolor{codestring}{rgb}{0.6,0.0,0.6}
\lstset{
    backgroundcolor=\color{codegray},
    basicstyle=\ttfamily\footnotesize,
    breaklines=true,
    captionpos=b,
    commentstyle=\color{codecomment},
    keywordstyle=\color{codekeyword}\bfseries,
    stringstyle=\color{codestring},
    language=Python,
    numbers=left,
    numberstyle=\tiny\color{gray},
    numbersep=5pt,
    showstringspaces=false,
    frame=single,
    tabsize=4,
    xleftmargin=15pt,
    framexleftmargin=10pt,
}

\title{Individual Report --- Student 4 \\
\large Ablation Studies, Visualization, and Theoretical Analysis \\
for the Dual-Arm Diffusion Policy Project}

\author{Student 4 \\
City University of Hong Kong}

\date{}

\begin{document}
\maketitle

% ============================================================
% 1. ROLE AND CONTRIBUTIONS
% ============================================================
\section{Role and Contributions}

In this team project, I served as Student 4, responsible for the
\textbf{full experiment pipeline, visualization, statistical analysis,
final report writing, and theoretical discussion}. My specific
deliverables include:

\begin{enumerate}
    \item \textbf{Ablation Experiments}: Designed and executed 12
    independent training runs of the Diffusion Policy under three
    controlled axes (demonstration quantity, prediction horizon, and
    observation horizon), totaling approximately 1.5 hours of GPU
    training on an NVIDIA RTX PRO 6000 (96GB).
    \item \textbf{Visualization Pipeline}: Implemented a unified
    matplotlib-based pipeline that renders all 12 figures in the team
    report, including success-rate comparisons, training-loss curves,
    t-SNE state distributions, and the radar plot summarizing
    cross-method trade-offs.
    \item \textbf{Statistical Analysis}: Aggregated raw outputs from
    Students 1--3 (BC test MSE, DAgger iteration logs, Diffusion Policy
    loss curves) into the final result tables and quantified covariate
    shift via Maximum Mean Discrepancy (MMD).
    \item \textbf{Final Report and Theoretical Discussion}: Authored
    the team's 12-page LaTeX report, including the theoretical
    treatment of covariate shift, the sim-to-real discussion, and the
    integration of all sub-team results into a coherent narrative.
\end{enumerate}

The work of other team members provided the inputs to my pipeline:
Student 1 produced the dual-arm Robosuite environment and 500 expert
demonstrations; Student 2 implemented BC and DAgger; Student 3 wrote
the Diffusion Policy training code (\texttt{ConditionalUnet1D}). My
ablation code reuses Student 3's model architecture without
modification, wrapping it in an automated experiment-management layer.

% ============================================================
% 2. ABLATION STUDIES
% ============================================================
\section{Ablation Studies}

\subsection{Motivation and Design}

Beyond reporting an aggregate success rate, the project specification
requires a controlled study of which design choices actually matter
for diffusion-based imitation learning. I therefore selected three
hyperparameters that span data, output structure, and input structure:

\begin{itemize}
    \item \textbf{Demonstration count} ($N \in \{50, 100, 200, 300, 500\}$):
    tests how data scale drives performance.
    \item \textbf{Prediction horizon} ($H_p \in \{4, 8, 16, 32\}$):
    tests how the length of the action chunk produced by the model
    affects learning.
    \item \textbf{Observation horizon} ($H_o \in \{1, 2, 4\}$):
    tests how much past context the model needs.
\end{itemize}

For each axis, only one variable changes at a time; all other
hyperparameters are fixed at the team default ($H_p = 16$, $H_o = 2$,
$N = 500$). All runs use \textbf{100 epochs} and
\textbf{batch size 256}; this is enough to observe converged trends
while keeping the total budget manageable, since I run 12 experiments
back-to-back. The model, optimizer (AdamW, $\text{lr} = 10^{-4}$,
$\text{wd} = 10^{-6}$), EMA (decay $0.999$), DDPM scheduler (100
steps, squared-cosine), and cosine LR schedule are identical to the
team default in Student 3's training code.

\subsection{Demonstration Quantity}

\begin{table}[H]
\centering
\small
\begin{tabular}{rrrr}
\toprule
\textbf{Demos} & \textbf{Samples} & \textbf{Final Loss} & \textbf{Time (min)} \\
\midrule
50  & 5{,}708  & 0.01875 & 1.2 \\
100 & 11{,}394 & 0.01053 & 2.0 \\
200 & 22{,}789 & 0.00625 & 3.9 \\
300 & 34{,}199 & 0.00393 & 5.8 \\
500 & 57{,}075 & 0.00220 & 9.6 \\
\bottomrule
\end{tabular}
\caption{Ablation on number of expert demonstrations.}
\label{tab:demo_count}
\end{table}

The dominant trend (Table~\ref{tab:demo_count}) is a near-linear
decrease in log-loss as demonstrations grow geometrically: each
doubling of $N$ reduces the final loss by roughly 40--45\%. This is
the strongest single effect observed in any of my ablations.
Importantly, the curve has not flattened at $N=500$, indicating that
the model is not yet data-saturated for this task complexity.

\paragraph{Practical implication.}
For diffusion-based imitation learning on coordinated dual-arm
manipulation, \emph{collecting more demonstrations gives consistently
larger returns than tuning the model architecture}. This finding
guides the team's recommendation that future work should prioritize
expanding the demonstration set rather than searching over model
hyperparameters.

\subsection{Prediction Horizon (Action Chunking)}

\begin{table}[H]
\centering
\small
\begin{tabular}{rrrr}
\toprule
\textbf{Pred Horizon} & \textbf{Samples} & \textbf{Final Loss} & \textbf{Time (min)} \\
\midrule
4  & 63{,}075 & 0.00207 & 10.6 \\
8  & 61{,}075 & 0.00251 & 10.3 \\
16 & 57{,}075 & 0.00224 & 9.7 \\
32 & 49{,}075 & 0.00229 & 8.5 \\
\bottomrule
\end{tabular}
\caption{Ablation on prediction horizon (action chunk length).}
\label{tab:pred_horizon}
\end{table}

Final losses lie within a narrow band of $0.00207$--$0.00251$
(Table~\ref{tab:pred_horizon}), even when the chunk length varies by
$8\times$. Two design choices explain this robustness:

\begin{enumerate}
    \item \textbf{Receding-horizon execution.} Although the model
    predicts $H_p$ steps, only the first $H_a = 2$ steps are executed
    before re-planning. The remaining steps act as a soft temporal
    prior, not a hard commitment, so high accuracy is required only
    over a short prefix.
    \item \textbf{Multi-scale 1D U-Net.} Longer sequences are
    absorbed by additional downsampling rather than by more
    learnable parameters per timestep, so the per-step prediction
    difficulty is roughly invariant to $H_p$.
\end{enumerate}

\paragraph{Caveat.} Training loss alone does not capture downstream
benefits of longer chunks (e.g., trajectory smoothness in closed-loop
execution). A definitive ablation would require online evaluation,
which we deferred due to environment access constraints.

\subsection{Observation Horizon (History Length)}

\begin{table}[H]
\centering
\small
\begin{tabular}{rrrr}
\toprule
\textbf{Obs Horizon} & \textbf{Cond Dim} & \textbf{Final Loss} & \textbf{Time (min)} \\
\midrule
1 & 12 & 0.00230 & 9.7 \\
2 & 24 & 0.00229 & 9.6 \\
4 & 48 & 0.00227 & 9.5 \\
\bottomrule
\end{tabular}
\caption{Ablation on observation horizon (history length).}
\label{tab:obs_horizon}
\end{table}

The three configurations differ by less than $1.5\%$ in final loss
(Table~\ref{tab:obs_horizon}). Physically, the pick-and-place task is
quasi-static at $20\,\text{Hz}$: the cube does not move on its own,
the end-effectors move at $\le 1\,\text{cm/step}$, and the goal is
fixed. A single observation frame already encodes the full Markov
state. Velocity or acceleration cues, which would require multiple
frames, contribute negligible information here.

\paragraph{When this finding would not generalize.} For tasks with
fast-moving objects (e.g., catching, hitting) or human-in-the-loop
intent estimation, multi-frame history is expected to matter much
more, since velocity is no longer derivable from a single frame.

\subsection{Cross-Ablation Summary}

Synthesizing the three studies:

\begin{itemize}
    \item \textbf{Data $\gg$ architecture.} The variance in final loss
    induced by the demo-count axis ($\approx 8.5\times$ between $N=50$
    and $N=500$) dwarfs the variance from either horizon axis
    ($\le 1.2\times$).
    \item \textbf{The default configuration is well-chosen.} The
    team's default ($H_p = 16$, $H_o = 2$) sits near the minimum of
    both horizon ablations, indicating that further hyperparameter
    search on these axes would yield diminishing returns.
    \item \textbf{The bottleneck is data, not capacity.} Since the
    model has not saturated at 500 demos, the practical next step is
    to enlarge the demonstration set, not to enlarge the network.
\end{itemize}

% ============================================================
% 3. VISUALIZATION AND STATISTICAL ANALYSIS
% ============================================================
\section{Visualization and Statistical Analysis}

\subsection{Unified Plotting Pipeline}

The script \texttt{student4\_visualizations.py} centralizes the
generation of all figures used in the team report. It enforces a
common visual style (color palette, font size, DPI, grid alpha) so
that figures sourced from disparate experiments remain visually
coherent. The pipeline produces:

\begin{itemize}
    \item Aggregate comparisons: success-rate bar chart
    (Fig.~1), method radar plot (Fig.~12), summary table image
    (Fig.~7).
    \item Training dynamics: BC training curve (Fig.~2), Diffusion
    training curve (Fig.~4), BC vs.\ Diffusion log-scale comparison
    (Fig.~5), four-phase Diffusion convergence breakdown (Fig.~9).
    \item Method-specific diagnostics: DAgger iteration improvement
    (Fig.~3), average episode length (Fig.~8).
    \item Conceptual figures: covariate-shift illustration
    (Fig.~10), action-chunking schematic (Fig.~11).
    \item Distribution analysis: t-SNE of states visited by Expert,
    BC, and DAgger policies (Fig.~6).
    \item Ablation outputs: training-loss curves and bar charts for
    all three ablation axes, plus a summary table image.
\end{itemize}

\subsection{Quantifying Covariate Shift}

To complement the qualitative t-SNE visualization, I computed the
Maximum Mean Discrepancy (MMD) between the state distributions of
the expert policy and each learned policy:

\begin{equation}
\text{MMD}^2(P, Q) = \mathbb{E}[k(x, x')] - 2\mathbb{E}[k(x, y)] + \mathbb{E}[k(y, y')]
\end{equation}

with an RBF kernel. The expert--BC MMD is substantially larger than
the expert--DAgger MMD, providing a quantitative confirmation of
what the t-SNE projection shows visually: BC visits states the expert
never reached, while DAgger's iterative on-policy correction pulls
its visitation distribution back toward the expert's.

% ============================================================
% 4. THEORETICAL DISCUSSION
% ============================================================
\section{Theoretical Discussion}

\subsection{Why BC Fails: A Compounding-Error View}

For a deterministic learned policy $\pi_\theta$ with single-step error
$\epsilon$ measured under the expert distribution, the worst-case
performance gap relative to the expert satisfies
\begin{equation}
J(\pi_\theta) - J(\pi^*) \leq T \cdot \epsilon + T^2 \cdot \epsilon^2,
\end{equation}
where $T$ is the task horizon. The quadratic term is the formal
statement of compounding error: at each step, the agent has some
probability of leaving the expert's support, and once it does, the
training signal provides no guidance. For our task ($T = 500$ steps),
even an offline MSE as small as $0.0058$ translates to a
$56$-percentage-point gap in success rate ($100\%$ vs.\ $44\%$).

\subsection{Why DAgger Helps}

DAgger replaces the worst-case quadratic dependence with an
expectation under $\pi_\theta$'s induced distribution, recovering an
$O(T \epsilon)$ regret bound. Empirically, our DAgger run rises from
$44\%$ to $86\%$ over its iterative training, with a non-monotonic
curve (drop at iteration 3, recovery thereafter). The non-monotonicity
is a known artifact of distribution shift in the aggregated dataset:
a temporarily poor policy contributes low-quality on-policy samples
that briefly degrade learning before later iterations recover.

\subsection{Why Diffusion Policy Is Different}

The Diffusion Policy attacks the same problem from a different angle:
rather than correcting the training distribution online, it models a
\emph{conditional distribution} $p(\mathbf{A}_t \mid \mathbf{O}_t)$
over action sequences. Two consequences follow.

\paragraph{Multi-modality.}
For states where multiple expert actions are valid (e.g., either arm
may grasp first), a deterministic regressor like BC averages the
modes and produces a physically inconsistent action. Diffusion Policy
samples one mode coherently.

\paragraph{Temporal commitment.}
Because each denoising pass produces an entire $H_p$-step chunk, the
sampled action sequence is internally consistent. Combined with
receding-horizon execution, this provides robustness to single-step
prediction noise that BC and DAgger cannot match.

In our experiments, the Diffusion Policy achieves a final training
loss of $4.4 \times 10^{-4}$, an order of magnitude lower than BC's
$5.8 \times 10^{-3}$. We estimate its online success rate at $\sim
90\%$; this estimate is necessarily imprecise, since training loss is
only a proxy for closed-loop performance, but the gap to BC's loss
strongly suggests Diffusion Policy would surpass DAgger if both were
evaluated under identical conditions.

% ============================================================
% 5. SIM-TO-REAL CONSIDERATIONS
% ============================================================
\section{Sim-to-Real Considerations}

The three methods carry different deployment implications:

\begin{itemize}
    \item \textbf{BC} is the cheapest to deploy but least robust to
    domain gaps. Any sim-to-real mismatch compounds through covariate
    shift.
    \item \textbf{DAgger} requires expert access during training. In
    simulation this is trivial (the scripted policy is always
    available); in the real world it usually means human teleoperation,
    which is expensive and slow. Variants such as HG-DAgger reduce the
    cost via intermittent corrections.
    \item \textbf{Diffusion Policy} is the most attractive for
    real-robot deployment: it trains entirely offline, naturally
    handles the multi-modal action distributions induced by domain
    randomization, and benefits from temporal robustness via action
    chunking. The main concern is inference latency from the 100-step
    denoising process; recent work on DDIM acceleration and
    consistency-distilled diffusion is directly applicable.
\end{itemize}

% ============================================================
% 6. LIMITATIONS
% ============================================================
\section{Limitations of My Contribution}

\begin{enumerate}
    \item \textbf{Training-loss-only ablation.} All ablation
    conclusions rest on training loss, not online success rate.
    Although loss is a strong correlate, online evaluation would be
    required for a definitive ranking, especially on the prediction-
    horizon axis where temporal smoothness is not captured by loss.
    \item \textbf{Single seed per configuration.} Each ablation cell
    is a single training run. The narrow gaps in the horizon
    ablations ($\le 1.5\%$) are within plausible seed-to-seed
    variance; replicate runs would tighten the conclusions.
    \item \textbf{Limited axes.} I varied three hyperparameters; the
    project specification also lists noise level as a candidate axis,
    which I did not explore.
    \item \textbf{Aggregate visualization, not raw re-evaluation.}
    Several figures (e.g., success-rate comparison) consume numbers
    produced by other team members. When their numbers update (as
    happened with the corrected DAgger result), the figures must be
    regenerated.
\end{enumerate}

% ============================================================
% 7. CONCLUSION
% ============================================================
\section{Conclusion}

My contribution to this project breaks down into four interlocking
deliverables: a controlled ablation study isolating the effect of
data quantity, prediction horizon, and observation horizon on
Diffusion Policy training; a unified visualization pipeline producing
all 12 figures used in the team report; a statistical synthesis of
the team's results into the main results table and the MMD
quantification of covariate shift; and the writing of the team's
final report including its theoretical and sim-to-real discussion.

The single most important quantitative finding from my ablations is
that \textbf{data quantity dominates hyperparameter choice} for this
task: an $8.5\times$ reduction in final loss from scaling demos $50
\to 500$, against $\le 1.2\times$ variation across either horizon
axis. This finding is operationally useful: it tells the team where
to invest effort in any follow-up work.

The single most important qualitative observation is that the three
methods do not occupy a single performance axis but instead trade off
along several dimensions --- success rate, training cost, sample
efficiency, multi-modal capability, and real-world deployability ---
which the team's radar plot makes explicit.

% ============================================================
% APPENDIX A: REPOSITORY
% ============================================================
\section*{Appendix A: Code Repository}

The complete source code, raw experiment outputs, and figures for my
contribution are released as a public GitHub repository:

\begin{center}
\fbox{\parbox{0.8\columnwidth}{\centering
\textbf{GitHub:} \url{https://github.com/Jacob0618/dual-arm-diffusion-ablation}
}}
\end{center}

The full ablation driver and visualization scripts are also reproduced
verbatim in Appendix C and Appendix D below.

% ============================================================
% APPENDIX B: SUBMISSION LAYOUT
% ============================================================
\section*{Appendix B: Submission Layout}

\begin{verbatim}
Student4_Submission/
  ablation_code/
    run_ablation.py          # 12-experiment driver (mine)
    model.py                 # ConditionalUnet1D (Student 3, unchanged)
  ablation_results/
    all_results.json         # final losses + configs
    demos_{50,100,200,300,500}/
    pred_horizon_{4,8,16,32}/
    obs_horizon_{1,2,4}/
    ablation_*.png           # rendered comparison plots
  visualization_code/
    student4_visualizations.py
  figures/                   # all 12 figures (PNG)
  figures_pdf/               # all figures (PDF, for LaTeX)
  report/
    student4_report.tex      # this document
  presentation.pptx
  presentation_v2.pptx
\end{verbatim}

% ============================================================
% APPENDIX C: ABLATION DRIVER CODE
% ============================================================
\newpage
\section*{Appendix C: Ablation Driver --- \texttt{run\_ablation.py}}

This is the complete script I wrote to drive all 12 ablation
experiments. The model architecture (\texttt{ConditionalUnet1D}) is
imported unmodified from Student 3's codebase.

\lstinputlisting[language=Python]{../ablation_code/run_ablation.py}

% ============================================================
% APPENDIX D: VISUALIZATION CODE
% ============================================================
\newpage
\section*{Appendix D: Visualization Pipeline --- \texttt{student4\_visualizations.py}}

This script renders all 12 figures used in the team report under a
unified style configuration.

\lstinputlisting[language=Python]{../visualization_code/student4_visualizations.py}

\end{document}
